\section{Introduction}



Two successful classes of probabilistic generative models involve sequentially corrupting training data with slowly increasing noise, and then learning to reverse this corruption in order to form a generative model of the data. 
\emph{Score matching with Langevin dynamics} (SMLD)~\citep{song2019generative} %
estimates the \emph{score} (\ie, the gradient of the log probability density with respect to data) at each noise scale, and then uses Langevin dynamics to sample from a sequence of decreasing noise scales during generation. 
\emph{Denoising diffusion probabilistic modeling} (DDPM)~\citep{sohl2015deep,ho2020denoising} trains a sequence of probabilistic models to reverse each step of the noise corruption, 
using knowledge of the functional form of the reverse distributions to make training tractable. 
For continuous state spaces, the DDPM training objective implicitly computes 
scores at each noise scale. 
We therefore refer to these two model classes together as \emph{score-based generative models}. %

Score-based generative models, and related techniques \citep{bordes2017learning,goyal2017variational, du2019implicit}, have proven effective at generation of images~\citep{song2019generative,song2020improved,ho2020denoising}, audio~\citep{chen2020wavegrad,kong2020diffwave}, graphs~\citep{niu20a}, and shapes~\citep{ShapeGF}. To enable new sampling methods and further extend the capabilities of score-based generative models, we propose a unified framework that generalizes previous approaches through the lens of stochastic differential equations (SDEs).%



\begin{SCfigure}
    \centering
    \caption{{\bf Solving a reverse-time SDE yields a score-based generative model.} Transforming data to a simple noise distribution can be accomplished with a continuous-time SDE. This SDE can be reversed if we know the score of the distribution at each intermediate time step, $\nabla_\bfx \log p_t(\bfx)$. }
    \includegraphics[width=0.65\linewidth, trim=105 0 110 0, clip]{figures/diffusion_schematic.pdf}
    \label{fig:teaser}
\end{SCfigure}

Specifically, instead of perturbing data with a finite number of noise distributions, we consider a continuum of distributions that evolve over time according to a diffusion process. This process progressively diffuses a data point into random noise, and is given by a prescribed SDE that does not depend on the data and has no trainable parameters. 
By reversing this process, we can smoothly mold random noise into data for sample generation. Crucially, this reverse process satisfies a reverse-time SDE~\citep{Anderson1982-ny}, which can be derived from the forward SDE given the score of the marginal probability densities as a function of time. We can therefore approximate the reverse-time SDE by training a time-dependent neural network to estimate the scores, and then produce samples using numerical SDE solvers. Our key idea is summarized in \cref{fig:teaser}.

Our proposed framework has several theoretical and practical contributions: %

\textbf{Flexible sampling and likelihood computation:}~ We can employ any general-purpose SDE solver to integrate the reverse-time SDE for sampling. In addition, we propose two special methods not viable for general SDEs: (i) Predictor-Corrector (PC) samplers that combine numerical SDE solvers with score-based MCMC approaches, such as Langevin MCMC~\citep{parisi1981correlation} and HMC~\citep{neal2011mcmc}; and (ii) deterministic samplers based on the probability flow ordinary differential equation (ODE). The former \emph{unifies and improves} over existing sampling methods for score-based models. The latter allows for \emph{fast adaptive sampling} via black-box ODE solvers, \emph{flexible data manipulation} via latent codes, a \emph{uniquely identifiable encoding}, and notably, \emph{exact likelihood computation}.

\textbf{Controllable generation:}~ We can modulate the generation process by conditioning on information not available during training, because the conditional reverse-time SDE can be efficiently estimated from \emph{unconditional} scores. This enables applications such as class-conditional generation, image inpainting, colorization and other inverse problems, all achievable using a single unconditional score-based model without re-training.


\textbf{Unified framework:}~ Our framework provides a unified way to explore and tune various SDEs for improving score-based generative models.
The methods of SMLD and DDPM can be amalgamated into our framework as discretizations of two separate SDEs. 
Although DDPM~\citep{ho2020denoising} was recently reported to achieve higher sample quality than SMLD~\citep{song2019generative,song2020improved}, we show that with better architectures and new sampling algorithms allowed by our framework, the latter can catch up---it achieves new state-of-the-art Inception score (9.89) and FID score (2.20) on CIFAR-10, as well as high-fidelity generation of $1024\times 1024$ images for the first time from a score-based model. In addition, we propose a new SDE under our framework that achieves a likelihood value of 2.99 bits/dim on uniformly dequantized CIFAR-10 images, setting a new record on this task.

